\documentclass[12pt,a4paper]{article}

\usepackage[italian]{babel}
\usepackage[T1]{fontenc}
\usepackage[utf8]{inputenc}
\usepackage{geometry}
\usepackage{graphicx}
\usepackage{listings}
\usepackage{lstautogobble}
\usepackage{xcolor}
\usepackage{hyperref}
\usepackage{booktabs}
\usepackage{amsmath}
\usepackage{amssymb}
\usepackage{float}
\usepackage{caption}
\usepackage{subcaption}
\usepackage{parskip}
\usepackage{microtype}
\usepackage{titlesec}
\usepackage{enumitem}
\usepackage{array}

\geometry{
  top=2.5cm,
  bottom=2.5cm,
  left=3cm,
  right=3cm
}

\hypersetup{
  colorlinks=true,
  linkcolor=black,
  urlcolor=blue,
  citecolor=black
}

% Stile per i listati di codice
\lstdefinestyle{codestyle}{
  backgroundcolor=\color{gray!8},
  basicstyle=\ttfamily\small,
  keywordstyle=\color{blue}\bfseries,
  commentstyle=\color{gray!60!black}\itshape,
  stringstyle=\color{green!50!black},
  numberstyle=\tiny\color{gray},
  breaklines=true,
  breakatwhitespace=false,
  numbers=left,
  numbersep=8pt,
  frame=single,
  framerule=0.4pt,
  rulecolor=\color{gray!40},
  captionpos=b,
  autogobble=true,
  showstringspaces=false,
  tabsize=4
}
\lstset{style=codestyle}

% Stile per JSON inline
\lstdefinelanguage{json}{
  morestring=[b]",
  literate=
    *{:}{{{\color{black}:}}}{1}
    {,}{{{\color{black},}}}{1}
    {\{}{{{\color{black}\{}}}{1}
    {\}}{{{\color{black}\}}}}{1}
    {[}{{{\color{black}[}}}{1}
    {]}{{{\color{black}]}}}{1},
  keywordstyle=\color{blue},
  stringstyle=\color{green!50!black},
  basicstyle=\ttfamily\small,
}

% Colonne tabelle
\newcolumntype{L}[1]{>{\raggedright\arraybackslash}p{#1}}
\newcolumntype{C}[1]{>{\centering\arraybackslash}p{#1}}
\newcolumntype{R}[1]{>{\raggedleft\arraybackslash}p{#1}}

% Titolo
\title{
  \vspace{-1cm}
  \Large\textbf{Integrazione di un Simulatore di Rete e di un Resource Manager}\\
  \large\textbf{in un Sistema di Percezione Cooperativa Multi-Agente}\\[0.5em]
  \normalsize Report tecnico del progetto di integrazione\\
  CoPerception -- OMNeT++ -- Konro
}
\author{}
\date{Aprile 2026}

\begin{document}

\maketitle
\tableofcontents
\newpage

% ============================================================
\section{Introduzione}
% ============================================================

Il presente documento descrive il lavoro svolto per estendere il framework di percezione cooperativa
\textit{CoPerception} con due componenti esterne: il simulatore di rete \textbf{OMNeT++} e il resource
manager \textbf{Konro}.

L'obiettivo complessivo del progetto è estendere il framework di percezione cooperativa multiagente con un simulatore di rete in grado di emulare più accuratamente l'applicazione del framework in contesti reali; dopodichè verificare se un resource manager basato su feedback applicativo possa migliorare il comportamento del sistema
in presenza di congestione o latenza elevata.

Il lavoro si articola nelle seguenti fasi principali:

\begin{enumerate}
  \item \textbf{Integrazione con OMNeT++}: il framework di deep learning viene collegato a un simulatore
        di rete realistico per sostituire la comunicazione ideale (istantanea) con una che tenga conto
        di latenza, perdite di pacchetti e posizione fisica degli agenti.

  \item \textbf{Integrazione con Konro}: il processo Python che esegue il modello invia periodicamente
        un segnale di feedback al resource manager, il quale pu\`o agire sulle risorse assegnate al
        processo in funzione della qualit\`a percepita.
\end{enumerate}

% ============================================================
\section{Il Framework CoPerception}
% ============================================================

CoPerception \`e un SDK Python/PyTorch per la percezione cooperativa in scenari di guida autonoma.
Fornisce un insieme di modelli, dataset e script di training/valutazione per task di rilevamento
oggetti (\textit{detection}) e segmentazione semantica da nuvole di punti 3D in formato BEV
(\textit{Bird's-Eye View}).

\subsection{Dataset V2X-Sim}

Il dataset utilizzato \`e \textbf{V2X-Sim}, un dataset sintetico che simula scenari di comunicazione
Vehicle-to-Everything (V2X) con fino a 6 agenti cooperanti per scena (5 veicoli e 1 RSU
opzionale). Ogni agente dispone di un sensore LiDAR e condivide le proprie feature con gli altri.
Il dataset \`e acquisito a \textbf{5 Hz} (un frame ogni 0,2 secondi).

\subsection{Pipeline di rilevamento cooperativo}

Per il task di detection il modello di riferimento \`e \textbf{DiscoNet} (Learning Distilled
Collaboration Graph for Multi-Agent Perception), che opera in modalit\`a di fusione intermedia
(\textit{intermediate fusion}): ogni agente calcola localmente le proprie feature map dalla
nuvola di punti LiDAR, le trasforma nel sistema di riferimento dell'agente ricevente tramite una
matrice di trasformazione omogenea $4\times4$, e le trasmette. Il ricevente fonde le feature
ricevute con le proprie prima dell'head di detection.

% ============================================================
\section{Integrazione con OMNeT++}
% ============================================================

\subsection{Motivazione}

Nel framework originale, la condivisione delle feature map tra agenti avviene in modo ideale:
i dati vengono copiati in memoria senza alcun ritardo e senza possibilit\`a di perdita.
Questo non riflette le condizioni reali di un canale wireless V2X (IEEE 802.11p), in cui
latenza, congestione e perdite di pacchetti sono eventi ordinari.

Per rendere la simulazione pi\`u fedele alla realt\`a, OMNeT++ viene usato come simulatore
di rete, capace di modellare la propagazione radio, la collisione di pacchetti e il
ritardo end-to-end in funzione della posizione fisica dei nodi.

\subsection{Architettura generale}

La connessione tra Python e OMNeT++ segue uno schema client-server TCP con un canale di
comunicazione basato su messaggi JSON delimitati da newline (\textit{newline-delimited JSON}).
(Figura~\ref{fig:arch}):

\begin{itemize}
  \item \textbf{OMNeT++} avvia un server TCP su una porta condivisa (default: 5555) attraverso
        il modulo \texttt{NetworkManager}.
  \item \textbf{Python} si connette al server all'inizio dell'esecuzione tramite la classe
        \texttt{OmnetBridge}.
\end{itemize}

\subsubsection{Paradigma asincrono event-based}

Una delle scelte di design fondamentali del sistema \`e l'adozione di un \textbf{paradigma
asincrono event-based} per lo scambio di informazioni tra i due processi. Questa scelta
\`e stata dettata da un vincolo strutturale preciso: i due sistemi operano su \textit{time
domain} incompatibili.

Il processo Python esegue inferenza frame per frame, dove ogni iterazione del loop di test
richiede un tempo di calcolo variabile e non deterministico (dipendente dal modello, dal
batch size, dalla GPU disponibile). OMNeT++, invece, avanza nel suo tempo simulato in passi
discreti, indipendentemente dall'orologio di sistema.

Obbligare i due processi a sincronizzarsi ad ogni scambio di pacchetto -- uno schema
\textit{request-response} sincrono -- avrebbe introdotto un accoppiamento temporale stretto
che avrebbe rallentato l'inferenza in attesa della risposta del simulatore per ogni coppia
mittente-ricevente. Con 6 agenti, ogni frame genera fino a $6 \times 5 = 30$ chiamate a
\texttt{feature\_transformation}: attendere una risposta sincrona per ciascuna avrebbe reso
il sistema inutilizzabile.

Il paradigma adottato \`e invece il seguente:

\begin{enumerate}
  \item Python invia a OMNeT++ i comandi \texttt{send} e \texttt{move} \textbf{in modo
        fire-and-forget}, senza attendere risposta.
  \item OMNeT++ processa i pacchetti nel proprio tempo simulato e, quando un pacchetto
        arriva a destinazione, invia un evento \texttt{received} di ritorno.
  \item Python legge gli eventi in arrivo da OMNeT++ \textbf{in modo non bloccante}
        (\texttt{update\_state()}) all'inizio di ogni chiamata alla monkey patch,
        aggiornando una cache locale (\texttt{latest\_delays}) con l'ultimo ritardo
        noto per ogni coppia mittente-ricevente.
  \item Il ritardo applicato alla feature map di un dato frame \`e l'ultimo valore
        ricevuto da OMNeT++ per quella coppia, non il valore del pacchetto corrente
        (che potrebbe non essere ancora stato consegnato).
\end{enumerate}

Questo schema introduce un \textbf{lag di un ciclo} nel feedback di ritardo: il ritardo
usato al frame $t$ \`e quello misurato da OMNeT++ per il pacchetto del frame $t-k$
(dove $k$ dipende dalla velocit\`a relativa dei due processi). Questo \`e accettabile
perch\'e i ritardi in un canale wireless reale variano lentamente rispetto alla frequenza
di campionamento del dataset (5 Hz), e la cache si aggiorna continuamente.

\subsubsection{Altre scelte di design}

\paragraph{Socket non bloccante lato Python.}
Il socket TCP aperto da \texttt{OmnetBridge} viene impostato immediatamente in modalit\`a
non bloccante dopo la connessione. Le letture (\texttt{recv}) vengono effettuate in loop
fino a \texttt{BlockingIOError}, che segnala che non ci sono pi\`u dati disponibili.
Questo evita che il processo Python si blocchi in attesa di un messaggio che potrebbe non
arrivare entro il frame corrente.

\paragraph{Handshake sincrono iniziale.}
L'unico punto sincrono dell'intero sistema \`e la fase di avvio: OMNeT++ attende in modo
bloccante la connessione di Python prima di avanzare nella simulazione. Questo garantisce
che i due processi siano allineati al momento zero e che non vengano persi comandi emessi
da Python prima che il server sia pronto. Una volta stabilita la connessione, il socket
viene immediatamente reso non bloccante e il loop di polling prende il controllo.

\paragraph{Separazione netta tra trasporto e decisione.}
Il modulo \texttt{CoPerceptionApp} in OMNeT++ si limita a inviare e ricevere pacchetti
UDP reali (con frammentazione) e a misurare il ritardo end-to-end. Non conosce la
semantica del payload (feature map, tensore PyTorch): per lui \`e semplicemente una
sequenza di byte. Tutta la logica applicativa (buffer, frame lag, sostituzione del
tensore) risiede nel wrapper Python. Questo disaccoppiamento mantiene il codice C++
semplice e riusabile indipendentemente dal modello di percezione utilizzato.

\paragraph{Fail-open per la robustezza.}
Se il bridge non riesce a connettersi a OMNeT++ o perde la connessione durante
l'esecuzione, il parametro \texttt{--network\_fail\_open} permette di continuare
l'inferenza degradando silenziosamente al comportamento ideale (ritardo zero, nessuna
perdita). Questo rende possibile eseguire esperimenti di baseline senza modificare
il codice, semplicemente non avviando OMNeT++.

\begin{figure}[H]
  \centering
  % Screenshot 1: schema architetturale dei due processi
  % Inserire qui un diagramma (anche disegnato a mano o con draw.io) che mostri:
  %   - il processo Python con OmnetBridge
  %   - il canale TCP sulla porta 5555
  %   - il processo OMNeT++ con NetworkManager e i nodi CoPerceptionApp
  \fbox{\parbox{0.8\textwidth}{\centering\vspace{2cm}
    [SCREENSHOT 1 -- Schema architetturale Python $\leftrightarrow$ OMNeT++]\\
    \small Catturare un diagramma o uno schema della connessione TCP tra i due processi.\\
    Si pu\`o realizzare con draw.io o analoga applicazione.
  \vspace{2cm}}}
  \caption{Schema dell'integrazione tra Python (CoPerception) e OMNeT++.}
  \label{fig:arch}
\end{figure}

\subsection{Il modulo NetworkManager}

Dal lato Omnet++, l'integrazione avviene tramite il modulo NetworkManager. La sua inizializzazione
segue due fasi distinte:

\begin{enumerate}
  \item \textbf{Handshake sincrono}: al lancio della simulazione, \texttt{NetworkManager} apre
        un socket TCP in modalit\`a bloccante e attende che Python si connetta. Questo garantisce
        che OMNeT++ non inizi a simulare prima che il modello di deep learning sia pronto.

  \item \textbf{Loop asincrono}: una volta stabilita la connessione, il socket viene impostato
        come non bloccante e viene avviato un ciclo di polling ogni 1 ms (\texttt{socketPoll})
        che legge e processa i comandi in arrivo.
\end{enumerate}

Il modulo registra tutti i sottomoduli \texttt{CoPerceptionApp} presenti nella simulazione
(uno per nodo veicolare) attraverso il metodo \texttt{discoverApps()}, che scansiona l'intera
gerarchia di moduli INET e mappa ogni agente al proprio indice.

\subsection{Il modulo CoPerceptionApp (lato OMNeT++)}

\texttt{CoPerceptionApp} \`e l'applicazione UDP che gira su ogni nodo veicolare simulato.
Le sue responsabilit\`a sono:

\begin{itemize}
  \item \textbf{Frammentazione}: le feature map hanno dimensioni dell'ordine del centinaio di
        kilobyte. Poich\'e il limite di UDP \`e 65\,535 byte, il payload viene spezzato in
        chunk da 60\,000 byte e inviato in frammenti successivi.
  \item \textbf{Misurazione del ritardo}: quando un pacchetto arriva a destinazione, l'app
        calcola il ritardo end-to-end come differenza tra il tempo di arrivo (\texttt{simTime()})
        e il tempo di creazione del pacchetto (\texttt{getCreationTime()}), e notifica il valore
        al \texttt{NetworkManager} tramite il metodo \texttt{notifyReception()}.
\end{itemize}

\subsection{Protocollo di comunicazione JSON}

Il protocollo prevede tre tipi di messaggi:

\subsubsection{Messaggio \texttt{move} (Python $\to$ OMNeT++)}

Aggiorna la posizione fisica di un nodo nel simulatore, estratta dalla matrice di trasformazione
omogenea del dataset. Questo permette al simulatore di calcolare la propagazione radio in funzione
della distanza reale tra agenti.

\begin{lstlisting}[language=json, caption={Messaggio move inviato da Python a OMNeT++ per aggiornare la posizione di un agente.}]
{"type": "move", "id": "2", "x": "15.3", "y": "-8.7", "z": "0.0"}
\end{lstlisting}

La posizione \`e espressa in metri nel sistema di riferimento dell'agente ego (agente 0).
L'estrazione avviene leggendo la colonna di traslazione dalla matrice di trasformazione
$\mathbf{T} \in \mathbb{R}^{4\times4}$:
\begin{equation}
  (x, y, z) = \bigl(T_{0,3},\; T_{1,3},\; T_{2,3}\bigr)
\end{equation}

\subsubsection{Messaggio \texttt{send} (Python $\to$ OMNeT++)}

Ordina al simulatore di inviare un pacchetto UDP dal nodo sorgente al nodo destinatario,
con la dimensione reale del tensore feature map in byte.

\begin{lstlisting}[language=json, caption={Messaggio send con dimensione reale del payload in byte.}]
{"type": "send", "src": "1", "dst": "0", "size": "98304", "id": "1->0"}
\end{lstlisting}

La dimensione viene calcolata come:
\begin{equation}
  \text{size\_bytes} = \text{element\_size} \times \prod_i \text{dim}_i
\end{equation}
dove \texttt{element\_size} \`e la dimensione in byte di un elemento del tensore PyTorch
(4 byte per \texttt{float32}).

\subsubsection{Messaggio \texttt{received} (OMNeT++ $\to$ Python)}

Notifica Python che un pacchetto \`e stato consegnato, insieme al ritardo misurato.

\begin{lstlisting}[language=json, caption={Risposta asincrona di OMNeT++ alla consegna di un pacchetto.}]
{"type": "received", "id": "1->0", "delay": 0.1823, "deliver": true}
\end{lstlisting}

\subsection{La Monkey Patch di \texttt{feature\_transformation}}

La tecnica usata per intercettare le comunicazioni tra agenti senza modificare il codice
del modello originale \`e la \textbf{monkey patch}: il metodo statico
\texttt{DetModelBase.feature\_transformation} viene sostituito a runtime con una funzione
\texttt{wrapped} che esegue la simulazione di rete prima di chiamare la funzione originale.

\begin{lstlisting}[language=python, caption={Sostituzione a runtime del metodo statico con la versione strumentata.}]
descriptor = DetModelBase.__dict__["feature_transformation"]
original = descriptor.__func__ if isinstance(descriptor, staticmethod) else descriptor

def wrapped(b, j, agent_idx, local_com_mat, all_warp, device, size, trans_matrices):
    # ... logica di rete ...
    local_com_mat_patched = local_com_mat.clone()
    local_com_mat_patched[b, j] = final_payload
    return original(b, j, agent_idx, local_com_mat_patched, all_warp, device, size, trans_matrices)

DetModelBase.feature_transformation = staticmethod(wrapped)
\end{lstlisting}

Al termine dell'esecuzione, una funzione \texttt{restore()} ripristina il metodo originale,
garantendo idempotenza anche in caso di errore (il blocco \texttt{finally} nella funzione
\texttt{main()} assicura sempre il ripristino).

\subsection{Gestione del buffer e del ritardo}

Per applicare correttamente il ritardo di rete sulle feature map, il wrapper mantiene un
\textbf{buffer circolare} per ciascun agente mittente, con capacit\`a massima di
\texttt{MAX\_BUFFER\_LEN = 10} frame (pari a 2 secondi al framerate di 5 Hz).

Quando OMNeT++ notifica un ritardo $d$ secondi, il \textit{frame lag} viene calcolato come:
\begin{equation}
  \text{frame\_lag} = \text{round}\!\left(d \times f_s\right)
\end{equation}
dove $f_s = 5\,\text{Hz}$ \`e il framerate del dataset.

Il comportamento in funzione del risultato della simulazione di rete \`e il seguente:

\begin{center}
\begin{tabular}{@{}lll@{}}
  \toprule
  Condizione & Azione & Log \\
  \midrule
  Pacchetto perso & Feature map sostituita con tensore di zeri & \texttt{[XXX]} \\
  Ritardo $< 0.1$\,s (lag = 0) & Feature map corrente utilizzata & \texttt{[OK]} \\
  Ritardo $\ge 0.1$\,s (lag $> 0$) & Feature map storica recuperata dal buffer & \texttt{[OLD]} \\
  Buffer insufficiente & Feature map pi\`u vecchia disponibile (underflow) & \texttt{[OLD!]} \\
  \bottomrule
\end{tabular}
\end{center}

\subsection{Configurazione OMNeT++}

La rete simulata comprende 6 nodi (\texttt{numHosts = 6}) con interfacce wireless IEEE 802.11p,
potenza di trasmissione 20\,mW e sensibilit\`a del ricevitore $-85$\,dBm. La mobilit\`a dei
nodi \`e gestita dal modulo personalizzato \texttt{PythonMobility}, che riceve le posizioni
aggiornate dal processo Python tramite i messaggi \texttt{move}.

\begin{figure}[H]
  \centering
  \fbox{\parbox{0.8\textwidth}{\centering\vspace{2cm}
    [SCREENSHOT 2 -- Terminale OMNeT++ in avvio]\\
    \small Catturare l'output del terminale di OMNeT++ che mostra:\\
    \texttt{NetworkManager: Listening on port 5555}\\
    \texttt{NetworkManager: Python connected! Starting simulation.}\\
    e i primi messaggi \texttt{move} e \texttt{send} processati.
  \vspace{2cm}}}
  \caption{Output di avvio del simulatore OMNeT++ con la fase di handshake TCP.}
  \label{fig:omnet_startup}
\end{figure}

\begin{figure}[H]
  \centering
  \fbox{\parbox{0.8\textwidth}{\centering\vspace{2cm}
    [SCREENSHOT 3 -- Log Python della monkey patch in esecuzione]\\
    \small Catturare alcune righe dell'output Python che mostrano i messaggi:\\
    \texttt{[OK]  1 -> 0 | Delay: 0.182s (Live) | phase=steady}\\
    \texttt{[OLD] 2 -> 0 | Delay: 0.400s -> Lag: 2 frames | phase=steady}\\
    \texttt{[XXX] 3 -> 0 | PACKET LOST | phase=steady}\\
    e le righe \texttt{[Proxy]} con i valori della metrica per frame.
  \vspace{2cm}}}
  \caption{Output della monkey patch durante un'esecuzione con OMNeT++ abilitato.}
  \label{fig:monkey_log}
\end{figure}

% ============================================================
\section{Integrazione con Konro}
% ============================================================

\subsection{Il resource manager Konro}

\textbf{Konro} \`e un resource manager progettato per gestire le risorse di calcolo di un
sistema multi-applicazione in funzione del feedback inviato dalle applicazioni stesse.
L'interfaccia esposta da Konro prevede due operazioni principali, accessibili via HTTP POST:

\begin{itemize}
  \item \texttt{/add}: registra il processo con Konro, comunicando PID, namespace e nome
        dell'applicazione.
  \item \texttt{/feedback}: invia periodicamente un valore intero compreso tra 0 e 200
        che rappresenta il rapporto tra il valore corrente della metrica e il target.
\end{itemize}

Il calcolo del feedback segue la formula:
\begin{equation}
  \text{feedback} = \text{clamp}\!\left(\left\lfloor\frac{v_\text{curr}}{v_\text{target}} \times 100\right\rfloor,\; 0,\; 200\right)
\end{equation}

Un valore di 100 indica che il target \`e esattamente raggiunto. Valori superiori a 100
segnalano che l'applicazione sta operando oltre il target e pu\`o essere ulteriormente
vincolata nelle risorse. Valori inferiori segnalano degrado della qualit\`a.

\subsection{Il bridge Python (konro\_bridge.py)}

Poich\'e il processo Python non pu\`o linkare direttamente la libreria C++ di Konro
senza modificare il sistema di build del progetto, \`e stato implementato un client HTTP
puro Python in \texttt{tools/det/konro\_bridge.py} che replica l'API di \texttt{konrolib}
usando la libreria \texttt{requests} (con fallback a \texttt{urllib}).

L'indirizzo del server Konro viene letto dalla variabile d'ambiente \texttt{KONRO}
(default: \texttt{http://localhost:8080}), in analogia con il comportamento della
libreria C++.

\subsection{La metrica proxy e la policy}

Per inviare un feedback significativo a Konro \`e necessario scegliere una metrica che
rifletta la qualit\`a percettiva del sistema in modo calcolabile frame per frame, senza
richiedere la valutazione completa della mAP (che richiederebbe l'accumulazione su
molti frame).

La metrica scelta \`e la seguente \textbf{proxy di qualit\`a per frame}:
\begin{equation}
  \text{proxy} = \text{recall} \times \min\!\left(1,\; \frac{N_\text{gt}}{N_\text{det}}\right)
\end{equation}
dove:
\begin{itemize}
  \item $\text{recall} = N_\text{tp} / \max(N_\text{gt}, 1)$ \`e la recall per frame,
        con $N_\text{tp}$ contato a IoU $\ge 0.5$.
  \item $N_\text{gt}$ \`e il numero di ground truth boxes nel frame corrente.
  \item $N_\text{det}$ \`e il numero di detections prodotte dal modello.
  \item Il fattore $\min(1, N_\text{gt}/N_\text{det})$ penalizza i falsi positivi:
        se il numero di detections supera significativamente il numero di ground truth,
        la qualit\`a scende.
\end{itemize}

La scelta di questa metrica \`e motivata dal fatto che:
\begin{enumerate}
  \item \`E calcolabile istantaneamente a ogni frame senza richiedere accumulo.
  \item Cattura sia i falsi negativi (bassa recall) sia i falsi positivi (penalizzazione).
  \item \`E normalizzata in $[0, 1]$ e quindi compatibile con il target di qualit\`a.
\end{enumerate}

\subsubsection{Smoothing con EMA}

La proxy grezzo per frame pu\`o oscillare notevolmente. Per ottenere un segnale di
controllo pi\`u stabile, viene applicato uno smoothing con \textbf{Exponential Moving
Average} (EMA) con fattore $\alpha = 0.2$:
\begin{equation}
  \text{ema}_t = \alpha \cdot \text{proxy}_t + (1-\alpha) \cdot \text{ema}_{t-1}
\end{equation}

Il valore EMA viene usato per calcolare il feedback inviato a Konro. Il target di
qualit\`a di default \`e $v_\text{target} = 0.85$. Il feedback viene inviato ogni frame
(intervallo configurabile).

\subsubsection{Monkey patch di cal\_local\_mAP}

Per intercettare le statistiche per frame senza modificare il test loop originale, viene
applicata una seconda monkey patch alla funzione \texttt{cal\_local\_mAP} di
\texttt{test\_codet}. La versione wrappata:
\begin{enumerate}
  \item Identifica quale agente sta venendo valutato in base al contatore di chiamate.
  \item Per l'agente monitorato (default: agente 1), chiama \texttt{cal\_frame\_stats}
        per ottenere $N_\text{gt}$, $N_\text{det}$, $N_\text{tp}$ e aggiorna il tracker.
  \item Chiama comunque la funzione originale per non interrompere l'accumulo della mAP.
\end{enumerate}

\subsection{Metriche misurate e loro significato}

Al termine di ogni esecuzione, il sistema produce un file JSON di riepilogo con due
categorie di metriche.

\subsubsection{Metriche di qualit\`a percettiva (sezione \texttt{proxy})}

\begin{center}
\begin{tabular}{@{}L{3.5cm}L{9cm}@{}}
  \toprule
  Metrica & Significato \\
  \midrule
  \texttt{proxy\_mean} & Media aritmetica della proxy su tutti i frame osservati.
                         Riflette la qualit\`a media del sistema. \\[4pt]
  \texttt{proxy\_ema} & Valore finale dell'EMA. Indica lo stato della qualit\`a
                        al termine della run. \\[4pt]
  \texttt{below\_target\_ratio} & Frazione di frame in cui la proxy \`e rimasta sotto
                                   il target 0.85. Un valore basso \`e desiderabile. \\
  \bottomrule
\end{tabular}
\end{center}

\subsubsection{Metriche di rete (sezione \texttt{network})}

\begin{center}
\begin{tabular}{@{}L{3.8cm}L{8.7cm}@{}}
  \toprule
  Metrica & Significato \\
  \midrule
  \texttt{delivery\_ratio} & Rapporto tra pacchetti consegnati e totale trasmessi. \\[4pt]
  \texttt{drop\_ratio} & Rapporto tra pacchetti persi e totale trasmessi. \\[4pt]
  \texttt{latency\_avg\_s} & Latenza media end-to-end in secondi per i pacchetti consegnati. \\[4pt]
  \texttt{latency\_p95\_s} & Latenza al 95-esimo percentile. Indica il worst case
                              quasi-estremo. \\[4pt]
  \texttt{stale\_packets} & Numero di pacchetti consegnati ma con lag $> 0$ frame
                            (dato obsoleto utilizzato). \\[4pt]
  \texttt{underflow\_packets} & Numero di volte in cui il buffer storico non era
                                 sufficiente a coprire il ritardo. \\
  \bottomrule
\end{tabular}
\end{center}

\subsection{Limiti della metrica proxy}

La metrica proxy presenta alcune limitazioni strutturali che \`e importante documentare.

\begin{enumerate}
  \item \textbf{Dipendenza dal ground truth}: il calcolo di $N_\text{gt}$ e $N_\text{tp}$
        richiede accesso alle annotazioni. In un sistema reale deployed, il ground truth
        non \`e disponibile. La proxy \`e quindi una metrica idealizzata, accessibile solo
        in fase sperimentale.

  \item \textbf{Monitoraggio di un solo agente}: il sistema attualmente monitora un solo
        agente alla volta (default: agente 1). L'esperienza di percezione degli altri
        agenti non contribuisce al feedback. Questo introduce un \textit{selection bias}
        nel segnale inviato a Konro.

  \item \textbf{Sensibilit\`a a scene sparse}: in frame con pochi o nessun oggetto
        ($N_\text{gt} = 0$, $N_\text{det} = 0$), il calcolo viene saltato per evitare
        divisioni per zero e spike artificiali della metrica.

  \item \textbf{Disaccoppiamento temporale}: il feedback inviato a Konro agisce sulle
        risorse del processo, ma l'effetto sulle richieste di rete avviene in modo
        indiretto. Konro non controlla direttamente la larghezza di banda del canale
        wireless, ma pu\`o influenzare il carico computazionale del nodo.

  \item \textbf{Lag del feedback}: l'EMA introduce un ritardo nel segnale di controllo.
        Con $\alpha = 0.2$, la costante di tempo effettiva \`e circa 4 frame (0.8 secondi
        al framerate di 5 Hz). Variazioni brusche della qualit\`a non vengono riflesse
        immediatamente nel feedback.
\end{enumerate}

% ============================================================
\section{Introduzione del Rumore Gaussiano sul Feedback}
% ============================================================

\subsection{Motivazione}

Lo studio sperimentale del rumore \`e di natura metodologica, non prestazionale.
L'obiettivo non \`e trovare un livello di rumore che migliori le prestazioni, ma
rispondere alla domanda: \textit{quanto deve essere precisa la metrica di feedback
affinch\'e la policy di Konro rimanga efficace?}

In un deployment reale, la metrica idealizzata basata sul ground truth non sarebbe
disponibile. Il sistema dovrebbe appoggiarsi a una metrica approssimata o appresa,
che necessariamente presenter\`a una deviazione dalla metrica ideale. Il rumore
gaussiano introdotto nei test modella proprio questa deviazione.

\subsection{Implementazione}

Il rumore viene aggiunto alla proxy grezza prima dell'aggiornamento dell'EMA,
con saturazione in $[0, 1]$:

\begin{lstlisting}[language=python, caption={Introduzione del rumore gaussiano sulla proxy di qualit\`a.}]
if self.feedback_noise_std > 0.0:
    proxy = min(1.0, max(0.0, proxy + random.gauss(0.0, self.feedback_noise_std)))
\end{lstlisting}

I due valori di deviazione standard testati sono:
\begin{itemize}
  \item $\sigma = 0.02$: piccola perturbazione, proxy molto vicina al valore ideale.
  \item $\sigma = 0.2$: perturbazione significativa, disaccoppiamento marcato tra
        proxy e qualit\`a reale.
\end{itemize}

\subsection{Risultati con OMNeT++ abilitato}

La Tabella~\ref{tab:omnet_enabled} riporta i valori delle metriche principali nei
quattro scenari con OMNeT++ abilitato.

\begin{table}[H]
  \centering
  \caption{Metriche per gli esperimenti con OMNeT++ abilitato.}
  \label{tab:omnet_enabled}
  \begin{tabular}{@{}L{4cm}R{2.2cm}R{2.2cm}R{2.2cm}R{2.2cm}@{}}
    \toprule
    Metrica & Baseline & Konro $\sigma=0$ & Konro $\sigma=0.02$ & Konro $\sigma=0.2$ \\
    \midrule
    \texttt{proxy\_mean}       & 0.6413 & 0.6615 & 0.6635 & 0.6374 \\
    \texttt{proxy\_ema}        & 0.8099 & 0.8378 & 0.8282 & 0.7222 \\
    \texttt{below\_target\_ratio} & 0.970 & 0.910 & 0.810 & 0.770 \\
    \texttt{delivery\_ratio}   & 0.947  & 0.947  & 0.947  & 0.947  \\
    \texttt{drop\_ratio}       & 0.053  & 0.053  & 0.053  & 0.053  \\
    \texttt{latency\_avg\_s}   & 0.1819 & 0.1587 & 0.1573 & 0.1526 \\
    \texttt{stale\_packets}    & 1241   & 557    & 552    & 515    \\
    \bottomrule
  \end{tabular}
\end{table}

\subsubsection{Analisi dei risultati}

\begin{itemize}
  \item \textbf{Baseline senza Konro}: \`e lo scenario peggiore in termini di
        latenza media e pacchetti obsoleti. Il sistema non si adatta alle condizioni
        di rete.

  \item \textbf{Konro, $\sigma = 0$ (feedback ideale)}: migliora tutte le metriche
        significative rispetto alla baseline. I pacchetti obsoleti si dimezzano
        (da 1241 a 557) e la latenza media scende da 0.182\,s a 0.159\,s.
        Il miglioramento non deriva da un diverso profilo di perdita (delivery ratio
        invariato) ma da una migliore adattabilit\`a del controller alle condizioni
        esistenti.

  \item \textbf{Konro, $\sigma = 0.02$}: \`e il risultato pi\`u forte dell'intero
        studio. Tutte le metriche migliorano rispetto alla baseline e non si
        registrano regressioni. Questo dimostra che un piccolo rumore sul feedback
        non compromette l'efficacia della policy.

  \item \textbf{Konro, $\sigma = 0.2$}: le metriche di rete continuano a migliorare
        rispetto alla baseline (latenza e pacchetti obsoleti ridotti ulteriormente),
        ma la qualit\`a proxy si degrada rispetto ai casi a rumore minore,
        specialmente nell'\texttt{ema} finale (0.722 vs 0.838). Il segnale di
        controllo diventa meno affidabile, ma il sistema non perde completamente
        i benefici dell'integrazione con Konro.
\end{itemize}

\begin{figure}[H]
  \centering
  \fbox{\parbox{0.8\textwidth}{\centering\vspace{2cm}
    [SCREENSHOT 4 -- Output del confronto tra run (compare\_konro\_runs.py)]\\
    \small Eseguire il confronto con:\\
    \texttt{python compare\_konro\_runs.py}\\
    e catturare la tabella prodotta che mostra le colonne\\
    \texttt{With Konro | Without Konro | Delta | Direction}\\
    per lo scenario OMNeT abilitato, indice 0 (sigma=0.0).
  \vspace{2cm}}}
  \caption{Tabella di confronto prodotta da \texttt{compare\_konro\_runs.py} per Konro senza rumore.}
  \label{fig:compare_table}
\end{figure}

\begin{figure}[H]
  \centering
  \fbox{\parbox{0.8\textwidth}{\centering\vspace{2cm}
    [SCREENSHOT 5 -- Output del confronto per sigma=0.02 e sigma=0.2]\\
    \small Ripetere la cattura per gli indici 1 (sigma=0.02) e 2 (sigma=0.2).\\
    Affiancare o impilare le due tabelle per confronto visivo.
  \vspace{2cm}}}
  \caption{Tabelle di confronto per le due configurazioni con rumore sul feedback.}
  \label{fig:compare_noise}
\end{figure}

\subsection{Risultati con OMNeT++ disabilitato}

Quando OMNeT++ \`e disabilitato, le metriche di rete sono identicamente zero
per costruzione: non c'\`e un canale fisico da simulare. Questo scenario funge da
\textbf{caso di controllo}: serve a verificare che il resource manager non produca
miglioramenti artificiali in assenza di un reale collo di bottiglia comunicativo.

\begin{table}[H]
  \centering
  \caption{Metriche per gli esperimenti con OMNeT++ disabilitato.}
  \label{tab:omnet_disabled}
  \begin{tabular}{@{}L{4cm}R{2.2cm}R{2.2cm}R{2.2cm}R{2.2cm}@{}}
    \toprule
    Metrica & Baseline & Konro $\sigma=0$ & Konro $\sigma=0.02$ & Konro $\sigma=0.2$ \\
    \midrule
    \texttt{proxy\_mean}       & 0.6775 & 0.6775 & 0.6757 & 0.6772 \\
    \texttt{proxy\_ema}        & 0.8403 & 0.8403 & 0.8416 & 0.7912 \\
    \texttt{below\_target\_ratio} & 0.900 & 0.900 & 0.820 & 0.700 \\
    \bottomrule
  \end{tabular}
\end{table}

\subsubsection{Analisi dei risultati}

\begin{itemize}
  \item Con $\sigma = 0$, il run Konro \`e effettivamente identico alla baseline.
        Il controller non ha nulla da ottimizzare perch\'e non esiste un canale
        di comunicazione reale.

  \item Con $\sigma = 0.02$, la \texttt{proxy\_mean} scende marginalmente ma
        il \texttt{below\_target\_ratio} migliora. L'effetto \`e trascurabile.

  \item Con $\sigma = 0.2$, l'\texttt{ema} scende pi\`u notevolmente (0.791).
        Questo conferma che un rumore elevato altera le dinamiche del controller
        anche in assenza di una rete attiva.
\end{itemize}

La conclusione di questo scenario \`e che il vantaggio reale di Konro emerge soltanto
quando il canale di comunicazione \`e stressato da OMNeT++.

\begin{figure}[H]
  \centering
  \fbox{\parbox{0.8\textwidth}{\centering\vspace{2cm}
    [SCREENSHOT 6 -- Riepilogo finale di un'esecuzione completa]\\
    \small Catturare l'output \texttt{=== Simulation Summary ===} al termine di una run\\
    con OMNeT abilitato e Konro abilitato (sigma=0.02), che mostra:\\
    \texttt{[Network] tx=... delivery=... drop=... avg=...s stale=...}\\
    \texttt{[Proxy] frames=... mean=... ema=... below\_target=...}
  \vspace{2cm}}}
  \caption{Output del riepilogo finale di una run completa.}
  \label{fig:summary}
\end{figure}

\begin{figure}[H]
  \centering
  \fbox{\parbox{0.8\textwidth}{\centering\vspace{2cm}
    [SCREENSHOT 7 -- Grafico di confronto (se disponibile)]\\
    \small Se il confronto produce un grafico PNG (\texttt{comparison\_latest.png}),\\
    inserirlo qui. In alternativa, catturare l'output grafico di matplotlib\\
    prodotto da \texttt{compare\_konro\_runs.py}.
  \vspace{2cm}}}
  \caption{Confronto grafico delle metriche chiave tra le diverse configurazioni di Konro.}
  \label{fig:chart}
\end{figure}

% ============================================================
\section{Riepilogo e Conclusioni}
% ============================================================

\subsection{Riepilogo dell'architettura integrata}

Il sistema finale integra tre componenti eterogenei:

\begin{enumerate}
  \item \textbf{CoPerception} (Python/PyTorch): esegue inferenza del modello DiscoNet
        sul dataset V2X-Sim, intercettata da due monkey patch che strumentano la
        comunicazione inter-agente e la valutazione per frame.

  \item \textbf{OMNeT++ con INET} (C++): simula il canale wireless 802.11p, calcola
        ritardi end-to-end realistici in funzione della posizione fisica dei nodi
        e notifica Python dei risultati di consegna.

  \item \textbf{Konro} (server esterno): riceve segnali di feedback sulla qualit\`a
        percettiva e agisce sulle risorse del processo; il suo comportamento \`e
        osservabile indirettamente attraverso l'evoluzione delle metriche di rete
        e di qualit\`a.
\end{enumerate}

\subsection{Conclusioni sull'efficacia di Konro}

I risultati mostrano chiaramente che:

\begin{itemize}
  \item Konro \`e efficace quando il sistema \`e esposto a condizioni di rete reali
        simulate da OMNeT++. I miglioramenti principali sono la riduzione dei
        pacchetti obsoleti e della latenza media.

  \item La policy \`e robusta a un livello moderato di rumore sul feedback
        ($\sigma = 0.02$): il sistema mantiene o migliora tutte le metriche rispetto
        alla baseline.

  \item Con un rumore pi\`u elevato ($\sigma = 0.2$), i benefici lato rete si
        conservano, ma la qualit\`a del segnale di controllo si degrada. Il sistema
        non fallisce, ma opera in modo meno ottimale.

  \item In assenza di OMNeT++, Konro non produce effetti significativi, confermando
        che il beneficio reale \`e legato alla presenza di un collo di bottiglia
        comunicativo.
\end{itemize}

\subsection{Sviluppi futuri}

\begin{itemize}
  \item Sostituire la proxy basata su ground truth con una metrica appresa o
        auto-supervisionata, per rendere il sistema deployable.
  \item Estendere il monitoraggio a tutti gli agenti contemporaneamente.
  \item Valutare l'impatto su metriche di detection pi\`u standard come mAP a
        diversi threshold IoU, accumulando risultati su pi\`u scene.
  \item Esplorare policy adattive che regolino anche la frequenza di invio delle
        feature map in funzione delle condizioni di rete.
\end{itemize}

% ============================================================
\appendix
\section{Guida agli screenshot}
% ============================================================

La Tabella~\ref{tab:screenshots} riassume tutti gli screenshot richiesti nel documento,
con la descrizione di dove e come catturarli.

\begin{table}[H]
  \centering
  \caption{Elenco degli screenshot da catturare.}
  \label{tab:screenshots}
  \small
  \begin{tabular}{@{}C{0.6cm}L{3.5cm}L{8.5cm}@{}}
    \toprule
    N. & Posizione nel testo & Cosa catturare \\
    \midrule
    1 & Figura~\ref{fig:arch} &
        Schema architetturale disegnato con draw.io o analogo, che mostri
        Python/OmnetBridge $\leftrightarrow$ TCP:5555 $\leftrightarrow$ OMNeT++/NetworkManager $\leftrightarrow$ CoPerceptionApp$\times$6. \\[6pt]
    2 & Figura~\ref{fig:omnet_startup} &
        Terminale OMNeT++ con la sequenza di avvio: \texttt{Listening on port 5555},
        \texttt{Python connected!}, i primi log \texttt{Registered app for vehicle ID}
        e i primi comandi \texttt{move}/\texttt{send} processati. \\[6pt]
    3 & Figura~\ref{fig:monkey_log} &
        Terminale Python durante l'inferenza: righe \texttt{[OK]}/\texttt{[OLD]}/\texttt{[XXX]}
        della monkey patch e righe \texttt{[Proxy]} con \texttt{recall}, \texttt{penalty},
        \texttt{proxy}, \texttt{ema} per frame. \\[6pt]
    4 & Figura~\ref{fig:compare_table} &
        Output testuale di \texttt{compare\_konro\_runs.py} per lo scenario OMNeT abilitato,
        indice 0 (senza rumore). Mostrare la tabella con le colonne
        \texttt{With Konro / Without Konro / Delta / Direction}. \\[6pt]
    5 & Figura~\ref{fig:compare_noise} &
        Stesso script per indici 1 ($\sigma=0.02$) e 2 ($\sigma=0.2$),
        affiancare o impilare. \\[6pt]
    6 & Figura~\ref{fig:summary} &
        Output \texttt{=== Simulation Summary ===} di una run completa con
        OMNeT e Konro abilitati ($\sigma=0.02$). \\[6pt]
    7 & Figura~\ref{fig:chart} &
        File \texttt{logs/ab/comparison\_latest.png} o grafico matplotlib prodotto
        da \texttt{compare\_konro\_runs.py}. \\
    \bottomrule
  \end{tabular}
\end{table}

\end{document}